% Reviewed translation: PDF pages 283--287 / printed pages 269--273.
% The original scan is authoritative; mathematical tokens were checked against it.
\MathBookSourcePage{283}
下面在 $h\to0$ 时正则剖分族上建立 $M_h$ 的逼近性质.

\begin{Theorem}\label{thm:incompressible-nedelec-approximation}
令 $\mathcal T_h$ 是 $\overline\Omega$ 的正则剖分族,
$M_h,r_h$ 由
\eqref{eq:incompressible-nedelec-global-space} 和
\eqref{eq:incompressible-nedelec-global-interpolant}
定义, $l\geq1$. 对所有 $\mathbf u\in H^{l+1}(\Omega)^3$,
\begin{equation}\label{eq:incompressible-nedelec-hcurl-approximation}
\|\mathbf u-r_h\mathbf u\|_{H(\operatorname{curl};\Omega)}
\leq C_1h^l\{\|\mathbf u\|_{l,\Omega}
+|\mathbf u|_{l+1,\Omega}\}.
\end{equation}
此外, 对 $s>2$ 和 $\mathbf u\in W^{1,s}(\Omega)^3$,
\begin{equation}\label{eq:incompressible-nedelec-stability}
\|\mathbf u-r_h\mathbf u\|_{0,\Omega}
+h\|\operatorname{curl}(\mathbf u-r_h\mathbf u)\|_{0,\Omega}
\leq C_2h|\mathbf u|_{1,s,\Omega},
\end{equation}
其中 $C_1,C_2$ 与 $h,\mathbf u$ 无关.
\end{Theorem}

\begin{Proof}
先证明 \eqref{eq:incompressible-nedelec-hcurl-approximation}.
由协变变换,
\[
\|\mathbf u-r_h\mathbf u\|_{0,\kappa}
\leq|\det B_\kappa|^{1/2}\|B_\kappa^{-1}\|
\|\widehat{\mathbf u}-r_{\widehat\kappa}
\widehat{\mathbf u}\|_{0,\widehat\kappa}.
\]
因参考插值算子保持 $P_{l-1}^3$, 附录插值推论给出
\[
\|\widehat{\mathbf u}-r_{\widehat\kappa}
\widehat{\mathbf u}\|_{0,\widehat\kappa}
\leq C_1
\begin{cases}
|\widehat{\mathbf u}|_{l,\widehat\kappa},&l\geq2,\\
|\widehat{\mathbf u}|_{1,\widehat\kappa}
+|\widehat{\mathbf u}|_{2,\widehat\kappa},&l=1.
\end{cases}
\]
又
\begin{equation}\label{eq:incompressible-nedelec-transformed-seminorm}
|\widehat{\mathbf u}|_{k,\widehat\kappa}
\leq C_2\|B_\kappa\|^{k+1}|\det B_\kappa|^{-1/2}
|\mathbf u|_{k,\kappa}.
\end{equation}
所以
\begin{equation}\label{eq:incompressible-nedelec-local-l2-estimate}
\begin{cases}
\|\mathbf u-r_h\mathbf u\|_{0,\kappa}
\leq C_3\|B_\kappa^{-1}\|\|B_\kappa\|^{l+1}
|\mathbf u|_{l,\kappa},&l\geq2,\\
\|\mathbf u-r_h\mathbf u\|_{0,\kappa}
\leq C_3\|B_\kappa^{-1}\|\|B_\kappa\|^2
(|\mathbf u|_{1,\kappa}+\|B_\kappa\||\mathbf u|_{2,\kappa}),&l=1.
\end{cases}
\end{equation}
由旋度变换,
\[
\|\operatorname{curl}(\mathbf u-r_h\mathbf u)\|_{0,\kappa}
\leq C_4|\det B_\kappa|^{1/2}\|B_\kappa^{-1}\|^2
\|\operatorname{curl}(\widehat{\mathbf u}
-r_{\widehat\kappa}\widehat{\mathbf u})\|_{0,\widehat\kappa}.
\]
Remark~\ref{rem:incompressible-nedelec-interpolation-curl-zero}
表明该线性映射在 $P_l^3$ 上为零, 故
\[
\|\operatorname{curl}(\widehat{\mathbf u}
-r_{\widehat\kappa}\widehat{\mathbf u})\|_{0,\widehat\kappa}
\leq C_5|\widehat{\mathbf u}|_{l+1,\widehat\kappa},
\]
进而
\begin{equation}\label{eq:incompressible-nedelec-local-curl-estimate}
\|\operatorname{curl}(\mathbf u-r_h\mathbf u)\|_{0,\kappa}
\leq C_6\|B_\kappa^{-1}\|^2\|B_\kappa\|^{l+2}
|\mathbf u|_{l+1,\kappa}.
\end{equation}
结合剖分正则性即得
\eqref{eq:incompressible-nedelec-hcurl-approximation}.

\MathBookSourcePage{284}
为证明稳定估计, 注意参考插值保持常数且
$r_{\widehat\kappa}\in
\mathcal L(W^{1,s}(\widehat\kappa)^3;R_l)$, 因而
\[
\|\widehat{\mathbf u}-r_{\widehat\kappa}
\widehat{\mathbf u}\|_{0,\widehat\kappa}
\leq C_7|\widehat{\mathbf u}|_{1,s,\widehat\kappa}.
\]
于是
\[
\|\mathbf u-r_h\mathbf u\|_{0,\kappa}
\leq C_8(\operatorname{meas}\kappa)^{1/2-1/s}
\|B_\kappa^{-1}\|\|B_\kappa\|^2|\mathbf u|_{1,s,\kappa}.
\]
Hölder 不等式和剖分正则性给出
\[
\|\mathbf u-r_h\mathbf u\|_{0,\Omega}
\leq C_9h(\operatorname{meas}\Omega)^{1/2-1/s}
|\mathbf u|_{1,s,\Omega}.
\]
同理,
$\|\operatorname{curl}(\widehat{\mathbf u}
-r_{\widehat\kappa}\widehat{\mathbf u})\|_{0,\widehat\kappa}
\leq C_{10}|\widehat{\mathbf u}|_{1,s,\widehat\kappa}$,
故
\[
\|\operatorname{curl}(\mathbf u-r_h\mathbf u)\|_{0,\Omega}
\leq C_{11}(\operatorname{meas}\Omega)^{1/2-1/s}
|\mathbf u|_{1,s,\Omega}.
\]
这完成证明.
\end{Proof}

\subsection{$l$ 次有限元的误差分析}
\label{subsec:incompressible-vector-potential-error-analysis}

$M_h,\Phi_h$ 已定义. 取标准有限元空间
\begin{equation}\label{eq:incompressible-vector-potential-scalar-space}
\Theta_h=\{q_h\in C^0(\overline\Omega);
\ q_h|_\kappa\in P_l\ \forall\kappa\in\mathcal T_h,
\ q_h|_\Gamma=0\}.
\end{equation}
回顾 $\Phi_{h0}$ 中函数满足
\eqref{eq:incompressible-vector-potential-discrete-divergence-free}.
下面验证离散 Poincaré 不等式和 curl lifting.

\begin{Lemma}\label{lem:incompressible-nedelec-gradient-interpolation}
若 $\mathbf u=\operatorname{grad}p$, $p\in H_0^1(\Omega)$,
且 $r_h\mathbf u$ 有定义, 则存在 $p_h\in\Theta_h$ 使
$r_h\mathbf u=\operatorname{grad}p_h$.
\end{Lemma}

\begin{Proof}
因为 $\operatorname{curl}\mathbf u=0$,
Remark~\ref{rem:incompressible-nedelec-interpolation-curl-zero}
给出 $\operatorname{curl}r_h\mathbf u=0$ 于每个单元.
又因 $p$ 在 $\Gamma$ 上为常数, 有
$\mathbf u\times\mathbf n=0$ 于 $\Gamma$.
Lemma~\ref{lem:incompressible-nedelec-conformity} 因而给出
$r_h\mathbf u\in H_0(\operatorname{curl};\Omega)$, 即
$\operatorname{curl}r_h\mathbf u=0$ 于 $\Omega$ 且
$r_h\mathbf u\times\mathbf n=0$ 于 $\Gamma$.

\MathBookSourcePage{285}
因此 $r_h\mathbf u=\operatorname{grad}q$
对某个 $q\in H_0^1(\Omega)$ 成立.
Lemma~\ref{lem:incompressible-nedelec-gradient-kernel}
又表明 $q|_\kappa\in P_l$, 故 $q\in\Theta_h$.
\end{Proof}

\begin{Remark}\label{rem:incompressible-nedelec-discrete-gradient-kernel}
对每个 $\boldsymbol\phi_h\in\Phi_h$,
若 $\operatorname{curl}\boldsymbol\phi_h=0$ 于 $\Omega$,
则存在唯一 $p_h\in\Theta_h$ 使
$\boldsymbol\phi_h=\operatorname{grad}p_h$. 因而
\[
\{\operatorname{grad}p_h;p_h\in\Theta_h\}
=\{\boldsymbol\phi_h\in\Phi_h;
\operatorname{curl}\boldsymbol\phi_h=0\}.
\]
\end{Remark}

为了使离散 Poincaré 不等式关于 $h$ 一致,
以下要求 $\mathcal T_h$ 一致正则.
\begin{Proposition}\label{prop:incompressible-nedelec-uniform-discrete-poincare}
设 $\Omega\subset\mathbb R^3$ 是具有多面体边界的有界凸开集,
$\mathcal T_h$ 是一致正则剖分. 则存在与 $h$ 无关的 $C^*$, 使
\begin{equation}\label{eq:incompressible-nedelec-uniform-discrete-poincare}
\|\boldsymbol\phi_h\|_{0,\Omega}
\leq C^*\|\operatorname{curl}\boldsymbol\phi_h\|_{0,\Omega}
\qquad\forall\boldsymbol\phi_h\in\Phi_{h0}.
\end{equation}
\end{Proposition}

\begin{Proof}
把 $\boldsymbol\phi_h$ 写成梯度与足够光滑的无散函数之和.
令 $p\in H_0^1(\Omega)$ 是
\[
(\operatorname{grad}p,\operatorname{grad}q)
=(\boldsymbol\phi_h,\operatorname{grad}q)
\qquad\forall q\in H_0^1(\Omega)
\]
的唯一解, 并置
$\mathbf w=\boldsymbol\phi_h-\operatorname{grad}p$.
则
\[
\operatorname{curl}\mathbf w
=\operatorname{curl}\boldsymbol\phi_h,
\quad\operatorname{div}\mathbf w=0,
\quad\mathbf w\times\mathbf n|_\Gamma=0.
\]
又 $\operatorname{curl}\boldsymbol\phi_h\in L^\gamma(\Omega)^3$
对所有 $\gamma$ 成立. 因 $\Omega$ 为凸域, 存在 $s>2$, 使
$\mathbf w\in W^{1,s}(\Omega)^3$ 且
\begin{equation}\label{eq:incompressible-nedelec-continuous-regularity}
\|\mathbf w\|_{1,\alpha,\Omega}
\leq C_1(\alpha)\|\operatorname{curl}\mathbf w\|_{0,\alpha,\Omega}
\qquad 2\leq\alpha\leq s.
\end{equation}
于是 $r_h\mathbf w$ 有定义. 因
$\boldsymbol\phi_h\in\Phi_h$,
$r_h(\operatorname{grad}p)$ 也有定义,
并由 Lemma~\ref{lem:incompressible-nedelec-gradient-interpolation}
存在 $p_h\in\Theta_h$ 使
$r_h(\operatorname{grad}p)=\operatorname{grad}p_h$.

\MathBookSourcePage{286}
从而
$\boldsymbol\phi_h=r_h\mathbf w+\operatorname{grad}p_h$.
在 \eqref{eq:incompressible-vector-potential-discrete-divergence-free}
中取 $q_h=p_h$, 得
$\|\boldsymbol\phi_h\|_{0,\Omega}leq\|r_h\mathbf w\|_{0,\Omega}$.
只需证明
\begin{equation}\label{eq:incompressible-nedelec-rh-w-bound}
\|r_h\mathbf w\|_{0,\Omega}
\leq C^*\|\operatorname{curl}\boldsymbol\phi_h\|_{0,\Omega}.
\end{equation}
由 \eqref{eq:incompressible-nedelec-stability} 和
\eqref{eq:incompressible-nedelec-continuous-regularity},
\[
\|\mathbf w-r_h\mathbf w\|_{0,\Omega}
\leq C_2h\|\operatorname{curl}\boldsymbol\phi_h\|_{0,s,\Omega}.
\]
一致正则性、旋度变换及逆估计给出
\[
\|\operatorname{curl}\boldsymbol\phi_h\|_{0,s,\Omega}
\leq C_3h^{3(1/s-1/2)}
\|\operatorname{curl}\boldsymbol\phi_h\|_{0,\Omega}.
\]
所以只要 $2<s\leq6$,
\[
\|\mathbf w-r_h\mathbf w\|_{0,\Omega}
\leq C_4h^\alpha
\|\operatorname{curl}\boldsymbol\phi_h\|_{0,\Omega}
\]
且 $\alpha\geq0$, 这证明
\eqref{eq:incompressible-nedelec-rh-w-bound}.
\end{Proof}

\begin{Corollary}\label{cor:incompressible-nedelec-discrete-decomposition}
令 $\Omega$ 是有界多面体. 对每个
$\mathbf w_h\in\Phi_h$, 存在唯一
$\mathbf v_h\in\Phi_{h0}$ 和 $p_h\in\Theta_h$, 使
\begin{equation}\label{eq:incompressible-nedelec-discrete-decomposition}
\mathbf w_h=\mathbf v_h+\operatorname{grad}p_h,
\qquad |p_h|_{1,\Omega}\leq\|\mathbf w_h\|_{0,\Omega}.
\end{equation}
在 Proposition~\ref{prop:incompressible-nedelec-uniform-discrete-poincare}
的假设下,
\begin{equation}\label{eq:incompressible-nedelec-discrete-divfree-bound}
\|\mathbf v_h\|_{H(\operatorname{curl};\Omega)}
\leq(1+C^{*2})^{1/2}
\|\operatorname{curl}\mathbf w_h\|_{0,\Omega}.
\end{equation}
\end{Corollary}

\begin{Proof}
取 $p_h\in\Theta_h$ 为
\[
(\operatorname{grad}p_h,\operatorname{grad}q_h)
=(\mathbf w_h,\operatorname{grad}q_h)
\qquad\forall q_h\in\Theta_h
\]
的唯一解. 则
$\mathbf v_h=\mathbf w_h-\operatorname{grad}p_h\in\Phi_{h0}$,
其余估计直接来自 Proposition~\ref{prop:incompressible-nedelec-uniform-discrete-poincare}.
\end{Proof}

Corollary~\ref{cor:incompressible-nedelec-discrete-decomposition}
的第一部分正好建立 curl lifting 条件.
结合本节所有结果得到:

\begin{Theorem}\label{thm:incompressible-vector-potential-final-error}
令 $\Omega\subset\mathbb R^3$ 为有界多面体.
由 \eqref{eq:incompressible-nedelec-global-space},
\eqref{eq:incompressible-nedelec-zero-trace-space} 和
\eqref{eq:incompressible-vector-potential-scalar-space}
定义的离散问题彼此等价, 且有唯一解
$u_h=(\operatorname{curl}\boldsymbol\psi_h,
\boldsymbol\omega_h)$.

若此外 $\Omega$ 为凸域, Stokes 问题
\eqref{eq:incompressible-vector-potential-stokes} 的解满足
\[
\boldsymbol\psi\in H^{l+1}(\Omega)^3,
\qquad
\boldsymbol\omega=-\Delta\boldsymbol\psi
\in H^{l+1}(\Omega)^3,
\qquad l\geq1,
\]
且 $\mathcal T_h$ 一致正则, 则
\begin{equation}\label{eq:incompressible-vector-potential-final-error}
\left\{
\begin{aligned}
\|\boldsymbol\omega-\boldsymbol\omega_h\|_{0,\Omega}
&\leq C_1\{h^{l-1}|\boldsymbol\psi|_{l+1,\Omega}
+h^l\|\boldsymbol\omega\|_{l+1,\Omega}\},\\
\|\operatorname{curl}(\boldsymbol\psi-\boldsymbol\psi_h)\|_{0,\Omega}
&\leq C_2\{(h^{l-1}+h^l)|\boldsymbol\psi|_{l+1,\Omega}
+h^l\|\boldsymbol\omega\|_{l+1,\Omega}\}.
\end{aligned}
\right.
\end{equation}
\end{Theorem}

\MathBookSourcePage{287}
\begin{Remark}\label{rem:incompressible-vector-potential-lost-order}
与二维情形相同, 下式导致损失一个 $h$ 的幂次:
\begin{equation}\label{eq:incompressible-vector-potential-consistency-loss}
\inf_{\boldsymbol\phi_h\in\Phi_{h0}}
\sup_{\boldsymbol\mu_h\in M_h}
\frac{|(\operatorname{curl}(\boldsymbol\psi-\boldsymbol\phi_h),
\operatorname{curl}\boldsymbol\mu_h)|}
{\|\boldsymbol\mu_h\|_{0,\Omega}}.
\end{equation}
若 $\Phi_0$ 上投影 $\mathring P_h\boldsymbol\psi\in\Phi_{h0}$,
\[
(\operatorname{curl}(\mathring P_h\boldsymbol\psi-
\boldsymbol\psi),\operatorname{curl}\boldsymbol\phi_h)=0
\quad\forall\boldsymbol\phi_h\in\Phi_{h0},
\]
满足对所有 $p\in[2,\infty]$, $s\in[1,l]$ 的估计
\begin{equation}\label{eq:incompressible-vector-potential-conjectured-lp-projection}
\|\mathring P_h\boldsymbol\psi-\boldsymbol\psi\|_{0,p,\Omega}
+h\|\operatorname{curl}(\mathring P_h\boldsymbol\psi-
\boldsymbol\psi)\|_{0,p,\Omega}
\leq Ch^{s+1}\|\boldsymbol\psi\|_{s+1,p,\Omega},
\end{equation}
则可用前述二维论证改进
\eqref{eq:incompressible-vector-potential-final-error},
并恢复部分丢失的 $h$ 次幂. 特别地, 一次元也会得到可接受的收敛率.
虽然 \eqref{eq:incompressible-vector-potential-conjectured-lp-projection}
仍是猜想, 但它并非不合理, 希望这一问题不久能得到解决.
\end{Remark}

\subsection{压力的间断逼近}
\label{subsec:incompressible-vector-potential-discontinuous-pressure}

本小节简述并分析求解
\eqref{eq:incompressible-vector-potential-constrained-discrete-problem}
和 \eqref{eq:incompressible-vector-potential-unconstrained-discrete-problem}
所对应压力项的有限元方法. 情形与
Subsection~\ref{subsec:incompressible-hhj-discontinuous-pressure}
非常相似, 所以几乎所有结果均不证自明地列出.
需构造 $H(\operatorname{div};\Omega)$ 的子空间 $D_h$,
使 $\boldsymbol\mu_h\in M_h$ 时
$\operatorname{curl}\boldsymbol\mu_h\in D_h$,
而 $\mathbf v_h\in D_h$ 时
$\operatorname{div}\mathbf v_h$ 属于离散压力空间.

\begin{Definition}\label{def:incompressible-vector-potential-pressure-reference-space}
对 $l\geq1$, 令
\[
D_l=P_{l-1}^3\oplus
\{p(\mathbf x)\mathbf x;
\ p\in\widetilde P_{l-1}\}.
\]
\end{Definition}
